\section{Conclusions}\label{sec:conclusion}

We investigated several established and modern adaptive thresholding techniques for the denoising and binarization of noisy color images, with a particular focus on efficient processing of PPM image data.
All methods were implemented according to their original formulations and evaluated with respect to both binarization quality and computational performance.

Our qualitative experiments show that classical approaches such as Niblack and Sauvola provide reasonable results but are less robust to background noise compared to more recent methods.
The Nick method improves upon these baselines, while the approaches of Bataineh and Bradley achieve the highest visual quality.
Bradley's method in particular produces the most consistent foreground-background separation across different window sizes.

In terms of efficiency, profiling-guided optimizations demonstrate that the computation of local window statistics is the dominant performance factor.
Integral image based implementations substantially reduce redundant calculations and provide significant runtime improvements, especially for high-resolution images.
Overall, our results highlight the importance of combining adaptive thresholding formulations with efficient implementation strategies.
